% Figure: psychrometric chart skeleton showing a cooling-and-dehumidification
% process. Axes are dry-bulb temperature (x) and humidity ratio w (y, g water
% per kg dry air). The saturation curve (phi = 100%) bounds the chart on the
% upper-left. A constant relative-humidity curve (phi = 50%) is drawn for
% reference. The process line runs from the warm humid supply state A
% (30 C, w = 0.0133) down and to the left to the cooled, dehumidified state B
% (14 C, w = 0.0090) on the coil. The path is shown as the engineering
% straight-line approximation A->B plus the true coil path that dips to the
% apparatus dew point. Saturation-curve points are steam-table w_sat values.
\begin{figure}[htbp]
\centering
\definecolor{engteal}{RGB}{30,120,120}
\begin{tikzpicture}
\begin{axis}[
  width=12cm, height=8cm,
  xlabel={dry-bulb temperature $T$ ($^{\circ}$C)},
  ylabel={humidity ratio $w$ (g water / kg dry air)},
  xmin=5, xmax=40, ymin=0, ymax=30,
  axis lines=left,
  tick label style={font=\scriptsize},
  label style={font=\small},
  legend pos=north west,
  legend style={font=\scriptsize},
  clip=true,
]
% saturation curve (phi = 100%): w_sat = 0.622 p_sat/(p - p_sat), in g/kg.
% steam-table p_sat (kPa) at T: 5->0.872, 10->1.228, 15->1.706, 20->2.339,
% 25->3.169, 30->4.246, 35->5.628, 40->7.384. Converted to w in g/kg at 101.325.
\addplot[engblue, very thick, smooth] coordinates {
  (5,5.4) (10,7.6) (15,10.6) (20,14.7) (25,20.0) (30,27.0)
};
\addlegendentry{saturation, $\phi=100\%$}
% phi = 50% reference curve: w at half the saturation partial pressure.
\addplot[enggray, thick, dashed, smooth] coordinates {
  (15,5.3) (20,7.3) (25,9.9) (30,13.3) (35,17.7) (40,23.3)
};
\addlegendentry{$\phi=50\%$}
% process: supply state A (30 C, w = 13.3 g/kg ~ phi 50%) to coil state B
% (14 C, w = 9.0 g/kg). Straight-line engineering approximation.
\addplot[engred, very thick] coordinates {(30,13.3) (14,9.0)};
\addlegendentry{cooling + dehumidification}
% true coil path dipping to the apparatus dew point near the saturation curve
\addplot[engorange, thick, dotted] coordinates {(30,13.3) (22,12.5) (16,10.4) (14,9.0)};
% state markers
\addplot[only marks, mark=*, engblack, mark size=2.0pt] coordinates {(30,13.3) (14,9.0)};
\node[engblack, font=\scriptsize, anchor=west] at (axis cs:30.4,13.3) {A (supply, $30\,^{\circ}$C, $50\%$)};
\node[engblack, font=\scriptsize, anchor=south west] at (axis cs:14.3,9.0) {B (coil exit)};
% dew-point dropline from A to the saturation curve
\addplot[enggray, dotted] coordinates {(30,13.3) (18.3,13.3)};
\node[enggray, font=\tiny, anchor=east] at (axis cs:18.0,13.3) {dew pt};
\end{axis}
\end{tikzpicture}
\caption[Cooling and dehumidification on a psychrometric chart]{A
cooling-and-dehumidification process drawn on the skeleton of a
psychrometric chart. The horizontal axis is dry-bulb temperature, the
vertical axis is the humidity ratio (mass of water vapour per mass of dry
air), and the upper-left boundary is the saturation curve $\phi=100\%$,
computed from the steam-table saturation pressure through
$w_{\text{sat}}=0.622\,p_{\text{sat}}/(p-p_{\text{sat}})$. Supply air enters
at state A ($30\,\degC$, $50\%$ relative humidity, $w=13.3\,\text{g}/
\text{kg}$). The air's dew point is the temperature where a horizontal line
from A meets the saturation curve, about $18\,\degC$. A cooling coil held
below this dew point first cools the air to saturation, then continues to
cool it along the saturation curve while liquid water drops out, leaving the
coil at state B ($14\,\degC$, $w=9.0\,\text{g}/\text{kg}$). The straight red
line is the engineering A-to-B approximation used for load accounting; the
dotted orange path is the true coil trajectory that dips toward the apparatus
dew point. The horizontal drop in $w$ from A to B is the moisture removed
(the latent load); the leftward drop in temperature is the sensible load.}
\label{fig:vol04ch02:psychrometric-process}
\end{figure}
